\saltoPag{}
\section{Anexo}
    \indent{Para el circuito utilizado anterior mente utilizar como entrada una función seno con 10$V_{pp}$ de amplitud y $50KHz$ de frecuencia.} 
    \sangria{} Para obtener la informacion sobre le osciloscopio y el generador de señales utilizados consultar la bibliografia al final del informe.
    \imagen{7cm}{./imagenes/señalGenerada.jpg}
    \sangria{} El osciloscopio estaba configurado con una escala de $5v$ por division. Por lo que se puede ver que la señal generada es de 10$V_{pp}$.
    \sangria{} Una vez generada la onda, se conecta al circuito y se mide la tensión de salida en la salida de cada resistencia.
    \begin{center} \textbf{---- Resistor $R_1$ ----} \end{center}
    \imagen{7cm}{./imagenes/señal10K.jpg}
    \sangria{} Para esta medicion se utilizo una escala de $2v$ por division, por lo que al tener 2 cuadrados y la mitad de una linea de amplitud(una linea equivale a $0,4v$), la tension de salida aproximada es de $8,4V_{pp}$.
    \begin{center} \textbf{---- Resistor $R_2$ ----} \end{center}
    \sangria{} Para esta medicion se utilizo una escala de $2v$ por division, por lo que al tener 2 cuadrados se puede decir que la tension de salida esta alrededor de $1,6V_{pp}$, considerando una amplitud $0,8v$.
    \imagen{7cm}{./imagenes/señal4,7K.jpg}
    \begin{center} \textbf{---- Resistor $R_2$ ----} \end{center}
    \sangria{} Para esta medicion se utilizo la misma configuracion que la anterior, por lo que la tension de salida aproximada tambien es de $1,6V_{pp}$, lo que tiene sentido ya que por los resistores en paralelo circula la misma tensión. 
